\newcommand{\figuraHidrometroComparacion}{%
\begin{figure}[H]
    \centering
    \begin{tikzpicture}[>=Latex,font=\small,line cap=round,line join=round]
        \foreach \x/\nivel/\titulo/\liq in {
            3.0/3.05/{(a) Agua, $S=1$}/{agua},
            9.0/3.05/{(b) Líquido más denso, $S_x>1$}/{líquido $x$}} {
            \draw[very thick] (\x-2.1,4.7)--(\x-2.1,0.35)--(\x+2.1,0.35)--(\x+2.1,4.7);
            \fill[blue!12] (\x-2.02,0.43) rectangle (\x+2.02,\nivel);
            \draw[very thick,blue!60!black] (\x-2.1,\nivel)--(\x+2.1,\nivel);
            \node[font=\bfseries] at (\x,5.05) {\titulo};
            \node[blue!60!black] at (\x-1.25,0.78) {\liq};
        }

        \foreach \x/\shift in {3.0/0,9.0/0.40} {
            \begin{scope}[shift={(0,\shift)}]
                \fill[gray!30] (\x,1.38) ellipse (0.58 and 0.88);
                \draw[very thick] (\x,1.38) ellipse (0.58 and 0.88);
                \fill[gray!55] (\x-0.38,0.88) rectangle (\x+0.38,1.18);
                \draw[very thick] (\x-0.13,2.08)--(\x-0.13,4.45)--(\x+0.13,4.45)--(\x+0.13,2.08);
                \foreach \y in {2.25,2.55,2.85,3.15,3.45,3.75,4.05}
                    \draw (\x+0.13,\y)--(\x+0.29,\y);
            \end{scope}
        }

        \draw[densely dashed,red!75!black] (2.55,3.05)--(3.62,3.05);
        \node[red!75!black,anchor=west] at (3.65,3.05) {$S=1.0$};
        \draw[densely dashed,red!75!black] (8.45,3.05)--(9.68,3.05);
        \node[red!75!black,anchor=west,font=\scriptsize] at (9.72,1.18)
            {menor inmersión};
        \node[align=center,text=black!70] at (6.0,0.15)
            {mayor densidad $\Rightarrow$ menor volumen sumergido};
    \end{tikzpicture}
    \caption{Comparación de un mismo hidrómetro en agua y en un líquido más
    denso. El equilibrio exige el mismo empuje en ambos casos; por ello el
    instrumento desplaza menos volumen y flota más alto en el líquido denso.}
    \label{fig:hidrometro-comparacion}
\end{figure}%
}